% !TEX root = ../Thesis.tex

%% Leitfragen für das Abstract:
%% - Motivation: Warum ist das Thema relevant?
%% - Fragestellung: Was untersucht die Arbeit?
%% - Methode: Wie wurde untersucht (Ansatz, Techniken)?
%% - Ergebnisse: Zu welchen Schlussfolgerungen kam die Forschung?
%% - Relevanz: Was trägt die Arbeit zum bestehenden Wissen bei?
%%
%% Format: 150-250 Wörter, prägnant und eigenständig verständlich

\chapter*{Kurzfassung}
\addcontentsline{toc}{chapter}{Kurzfassung}

% Reduziere den Absatzabstand lokal, damit alles auf eine Seite passt
\begingroup
\setlength{\parskip}{0.5em}
Die Open-Source-Plattform frag.jetzt ermöglicht anonyme Echtzeitfragen in der Hochschullehre. Ihre verteilte Architektur aus mehreren Backend-Diensten, einer relationalen Datenbank und einem Message Broker wurde bislang ohne aggregierte Telemetrieerfassung betrieben. Störungen wurden erst wahrgenommen, wenn Endnutzende direkt betroffen waren. Die vorliegende Arbeit entwickelt eine Observability-Strategie, die dieses Defizit adressiert. Ausgehend von einer Systemanalyse mit sechs identifizierten Risikopunkten und fünf Observability-Anforderungen wird eine servicebezogene Operationalisierung der Four Golden Signals (Latenz, Traffic, Fehler, Sättigung) vorgenommen. Eine kriteriengestützte Evaluation dreier Architekturkandidaten führt zur Auswahl des LGTM-Stacks (Loki, Grafana, Tempo, Mimir/Prometheus) mit OpenTelemetry als Instrumentierungsschicht. Der wesentliche Beitrag der Arbeit liegt in einer methodisch nachvollziehbaren Überführung allgemeiner Observability-Konzepte in eine konkrete, dienstbezogene Strategie für ein reales verteiltes System. Die konzipierte Strategie umfasst eine serviceübergreifende Signal-Matrix, ein rollenbasiertes Dashboard-Konzept, symptomorientierte Alerting-Regeln mit Multi-Signal-Logik sowie strukturierte Runbooks. Die prototypische Umsetzung auf dem Staging-Server von frag.jetzt validiert die Kernmechanismen: Distributed Tracing über den OTel Java Agent, bidirektionale Log-Trace-Korrelation und ein Golden-Signals-Dashboard. Die Evaluation zeigt, dass vier von fünf Anforderungen weitgehend oder vollständig erfüllt werden, während die Ende-zu-Ende-Nachvollziehbarkeit aufgrund unvollständiger Instrumentierungsabdeckung teilweise erfüllt bleibt. Die Arbeit liefert übertragbare methodische Artefakte für vergleichbare cloud-native Plattformen.

\textbf{Schlüsselwörter:} Observability, Four Golden Signals, OpenTelemetry, Distributed Tracing, Alerting, Grafana, Microservices
\par
\endgroup

% Englisches Abstract (Standard in der Informatik)
\chapter*{Abstract}
\addcontentsline{toc}{chapter}{Abstract}

\begingroup
\setlength{\parskip}{0.5em}
The open-source platform frag.jetzt enables anonymous real-time questions in higher education. Its distributed architecture, comprising multiple backend services, a relational database, and a message broker, had been operated without aggregated telemetry collection. Failures only became visible when end users were already affected. This thesis develops an observability strategy to address this deficit. Starting from a system analysis that identifies six risk points and five observability requirements, the Four Golden Signals (latency, traffic, errors, saturation) are operationalized on a per-service basis. A criteria-driven evaluation of three candidate architectures leads to the selection of the LGTM stack (Loki, Grafana, Tempo, Mimir/Prometheus) with OpenTelemetry as the instrumentation layer. The thesis’ main contribution lies in the methodologically transparent transfer of general observability concepts into a concrete service-specific strategy for a real-world distributed system. The resulting strategy comprises a cross-service signal matrix, a role-based dashboard concept, symptom-oriented alerting rules with multi-signal logic, and structured runbooks. A prototypical implementation on the frag.jetzt staging server validates the core mechanisms: distributed tracing via the OTel Java Agent, bidirectional log-trace correlation, and a Golden Signals dashboard. The evaluation demonstrates that four of five requirements are largely or fully met, while end-to-end traceability remains partially fulfilled due to incomplete instrumentation coverage. The thesis contributes transferable methodological artifacts applicable to comparable cloud-native platforms.

\textbf{Keywords:} Observability, Four Golden Signals, OpenTelemetry, Distributed Tracing, Alerting, Grafana, Microservices
\par
\endgroup

\cleardoublepage
